% !TEX TS-program = lualatex

\documentclass[border = 3pt]{standalone}

%%%
% Des couleurs faciles d'emploi via le package \xcolor.
%%%
\usepackage[svgnames]{xcolor}

%%%
% La bibliothèque \luadraw allie une facilité d’utilisation
% à un rendu particulièrement soigné.
%%%
\usepackage[3d]{luadraw}

\begin{document}

\begin{luadraw*}{name = stairway}
local ld = luadraw

local Origin = ld.pt3d.Origin
local M      = ld.pt3d.M

local vecI = ld.pt3d.vecI
local vecJ = ld.pt3d.vecJ
local vecK = ld.pt3d.vecK

------
-- Nombre de marches, ou de tranches, souhaité.
------
local nb_slices = 10

------
-- ¨Def de la zone graphique.
------
local graphview = ld.graph3d:new{
  window3d = {-5, 5, -5, 5, -5, 5},
  viewdir  = {-35, 45},
  adjust2d = true,
  bbox     = false
}

------
-- Tracé de l'escalier avec l'obligation de tracer les éléments
-- du "plus caché" au "moins caché" (¨luadraw n'est pas un moteur
-- de rendu 3D).
--
--
-- note::
--     En décommentant la ligne suivante, un repère en mode
--     "boîte" sera visible. Ceci permet de choisir le bon ordre
--     des tracés suivant les réglages de la caméra.
------
graphview:Dboxaxes3d()

local num_col = 0
local colors  = {"ForestGreen", "Crimson"}

local cube = ld.parallelep(
  M(-5, -5, -5),
  vecI, vecJ, vecK
)

nb_slices = nb_slices - 1

for i = 0, nb_slices do
  num_col = num_col % 2 + 1

-- Attention aux ordres des boucles !
  for j = nb_slices, 0, -1 do
    for k = 0, nb_slices - i do
      if j == 0 or k == nb_slices - i then
        graphview:Dpoly(
          ld.shift3d(
            cube,
            i*vecI + j*vecJ + k*vecK
          ),
------
-- note::
--     En utilisant ''mode = ld.mFlatHidden'' en plus de la
--     couleur, les traits cachés apparaitront en pointillés.
------
          {color = colors[num_col]}
        )
      end
    end
  end
end

------
-- Montrons le résultat de notre oeuvre.
------
graphview:Show()
\end{luadraw*}

\end{document}
